%! Unidad 9 — Cinemática aplicada: ejercicios resueltos, ejercicios finales, olimpiadas y autoevaluación
%! Fuente: Fuentes/Fisica/Unidad 9 («Ejercicios resueltos», «Practica y calcula», «Piensa y resuelve») y Unidad 10
%! (captura 217: «Olimpiadas Física» y «Autoevaluación» de este tema). Enunciados adaptados con datos cambiados (y más realistas
%! cuando los de la fuente no lo eran); soluciones rehechas con g = 9,8 m/s². Dificultad: \(\star\) baja, \(\star\star\) media,
%! \(\star\star\star\) alta.

\section{Ejercicios resueltos}\label{sec:ud9-resueltos}

\begin{example}
    \textbf{MRU y MRUA: una persecución.} Mientras estás en una tienda, una persona coge tu bicicleta y empieza a pedalear desde el
    reposo con una aceleración de \qty{1.5}{\m\per\s\squared}. Un coche de policía pasa por delante de la tienda \qty{4}{\s} después,
    a \qty{54}{\km\per\hour} constantes, en la misma dirección. Calcula: a) cuánto tarda la policía en alcanzar a la persona;
    b) a qué distancia de la tienda la alcanza.

    \textbf{Solución.} \textbf{1. Datos}: fugitivo, MRUA con \(v_0 = 0\) y \(a = \qty{1.5}{\m\per\s\squared}\); policía, MRU con
    \(v = \qty{15}{\m\per\s}\) que empieza en la tienda en \(t_0 = \qty{4}{\s}\). Origen en la tienda y \(t = 0\) cuando arranca la bici.

    \textbf{2. Ecuaciones}: \(x_F = \num{0.75}\,t^2\) y \(x_P = 15\,(t - 4)\).

    \textbf{3. a)} La alcanza cuando \(x_F = x_P\): \(\num{0.75}\,t^2 = 15t - 60 \Rightarrow t^2 - 20t + 80 = 0\), con soluciones
    \(t = \qty{5.53}{\s}\) y \(t = \qty{14.5}{\s}\). La primera es el instante en que la policía la alcanza (la segunda sería
    cuando la bicicleta, cada vez más rápida, volvería a adelantar al coche si este no se detuviera).
    \textbf{b)} \(x = \num{0.75} \cdot \num{5.53}^2 = \qty{22.9}{\m}\). Comprobación: \(x_P = 15 \cdot \num{1.53} = \qty{22.9}{\m}\).
\end{example}

\begin{example}
    \textbf{Tiro parabólico: un balón.} Una futbolista golpea un balón desde el suelo a \qty{72}{\km\per\hour} con un ángulo de
    \(35^\circ\) con la horizontal. Determina: a) la altura máxima; b) el tiempo que el balón está en el aire; c) su alcance.

    \textbf{Solución.} \textbf{1. Datos}: \(v_0 = \qty{20}{\m\per\s}\), \(\alpha = 35^\circ\): \(v_{0x} = 20\cos 35^\circ =
    \qty{16.4}{\m\per\s}\); \(v_{0y} = 20\,\sen 35^\circ = \qty{11.5}{\m\per\s}\). \textbf{Incógnitas}: \(y_{\text{máx}}\),
    \(t_{\text{vuelo}}\) y \(x_{\text{máx}}\).

    \textbf{2. a)} En el punto más alto, \(v_y = 0\): \(t = \dfrac{\num{11.47}}{\num{9.8}} = \qty{1.17}{\s}\);
    \(y_{\text{máx}} = \num{11.47} \cdot \num{1.17} - \num{4.9} \cdot \num{1.17}^2 = \qty{6.71}{\m}\).

    \textbf{3. b)} Como sale y llega al mismo nivel, tarda lo mismo en bajar que en subir: \(t_{\text{vuelo}} = \qty{2.34}{\s}\).
    (Comprobación: la caída libre desde \qty{6.71}{\m} dura \(\sqrt{\num{6.71}/\num{4.9}} = \qty{1.17}{\s}\).)

    \textbf{4. c)} \(x_{\text{máx}} = \num{16.38} \cdot \num{2.34} = \qty{38.4}{\m}\).
\end{example}

\begin{example}
    \textbf{MCUA: una rueda.} Una rueda empieza a girar con una aceleración angular de \qty{4}{\radian\per\s\squared} desde una
    posición angular inicial de \(45^\circ\). Calcula: a) las ecuaciones de la posición angular y de la velocidad angular; b) el
    desplazamiento angular y la velocidad angular \qty{6}{\s} después de empezar a girar.

    \textbf{Solución.} \textbf{1. Datos}: \(\alpha = \qty{4}{\radian\per\s\squared}\), \(\omega_0 = 0\), \(\varphi_0 = 45^\circ =
    \pi/4~\unit{\radian}\).

    \textbf{2. a)} \(\varphi = \dfrac{\pi}{4} + \dfrac{1}{2} \cdot 4\,t^2 = \dfrac{\pi}{4} + 2t^2\) (rad); \(\omega = 4t\) (rad/s).

    \textbf{3. b)} \(\Delta\varphi = \varphi - \varphi_0 = 2 \cdot 6^2 = \qty{72}{\radian} = \dfrac{72}{2\pi} = \num{11.5}\) vueltas;
    \(\omega(6) = \qty{24}{\radian\per\s}\).
\end{example}

\begin{example}
    \textbf{Tiro parabólico: una bomba volcánica.} Los volcanes en erupción lanzan fragmentos de lava incandescente, las bombas
    volcánicas. En la erupción de 2021 en La Palma, imagina que una bomba sale de la boca de un cono a \qty{1100}{\m} sobre el nivel
    del mar con un ángulo de \(50^\circ\) y destruye un muro situado a \qty{700}{\m} de distancia horizontal y a \qty{600}{\m} de
    altitud. Calcula: a) la velocidad de lanzamiento; b) el tiempo de vuelo hasta el muro.

    \textbf{Solución.} \textbf{1. Datos}: origen al nivel del mar, en la vertical de la boca: \(y_0 = \qty{1100}{\m}\); \(\alpha =
    50^\circ\); punto final \((700, 600)\). \textbf{Incógnitas}: \(v_0\) y \(t\).

    \textbf{2. Ecuaciones}: \(x = v_0\cos 50^\circ\,t\) e \(y = 1100 + v_0\,\sen 50^\circ\,t - \num{4.9}\,t^2\).

    \textbf{3.} Al llegar al muro: \(700 = v_0\cos 50^\circ\,t\) y \(600 = 1100 + v_0\,\sen 50^\circ\,t - \num{4.9}\,t^2\). De la
    primera, \(v_0\,t = \dfrac{700}{\cos 50^\circ}\), y entonces \(v_0\,\sen 50^\circ\,t = 700\,\tg 50^\circ = \qty{834}{\m}\).
    Sustituyendo en la segunda: \(600 = 1100 + 834 - \num{4.9}\,t^2\), de donde \(t = \sqrt{\dfrac{1334}{\num{4.9}}} = \qty{16.5}{\s}\).

    \textbf{4.} \(v_0 = \dfrac{700}{\cos 50^\circ \cdot \num{16.5}} = \qty{66.0}{\m\per\s}\), unos \qty{238}{\km\per\hour}.
\end{example}

\begin{example}
    \textbf{Tiro horizontal: ayuda humanitaria.} Un helicóptero que vuela horizontalmente a \qty{180}{\km\per\hour} y a
    \qty{2000}{\m} de altura suelta un paquete. Calcula: a) el tiempo que tarda el paquete en llegar al suelo; b) a qué distancia
    horizontal del objetivo debe soltarlo; c) dónde está el helicóptero cuando el paquete llega al suelo.

    \textbf{Solución.} El paquete sale con la velocidad del helicóptero, \(v_0 = \qty{50}{\m\per\s}\) horizontal:
    \(x = 50t\), \(y = 2000 - \num{4.9}\,t^2\).
    a) \(y = 0\): \(t = \sqrt{\dfrac{2000}{\num{4.9}}} = \qty{20.2}{\s}\). b) \(x = 50 \cdot \num{20.2} = \qty{1010}{\m}\) antes del
    objetivo. c) Como el paquete avanza en horizontal igual que el helicóptero, este está \textbf{justo en su vertical}, a
    \qty{2000}{\m} de altura (si no se tiene en cuenta el rozamiento con el aire).
\end{example}

%========================================================================================
\section{Ejercicios de aplicación}\label{sec:ud9-aplicacion}
%========================================================================================

\subsection{Movimiento rectilíneo uniforme}

\begin{exercise}
    \(\star\) La luz del Sol tarda 8 minutos y 20 segundos en llegar a la Tierra. ¿A qué distancia está el Sol? (Velocidad de la luz:
    \qty{3e8}{\m\per\s}.)
    \begin{solu}
        \(d = \num{3e8} \cdot 500 = \textcolor{red}{\qty{1.5e11}{\m}}\).
    \end{solu}
\end{exercise}

\begin{exercise}
    \(\star\star\) Una partícula se mueve en línea recta con velocidad constante. Pasa por la posición \qty{40}{\m} en
    \(t = \qty{2}{\s}\) y por la posición \qty{16}{\m} en \(t = \qty{5}{\s}\). a) Representa la gráfica \(x\)-\(t\). b) Calcula su
    velocidad y su posición inicial. c) ¿Qué espacio recorre en \qty{6}{\s}?
    \begin{solu}
        a) Una recta descendente. b) \(v = \dfrac{16 - 40}{5 - 2} = \textcolor{red}{\qty{-8}{\m\per\s}}\);
        \(x_0 = 40 + 8 \cdot 2 = \textcolor{red}{\qty{56}{\m}}\). c) \(\textcolor{red}{\qty{48}{\m}}\).
    \end{solu}
\end{exercise}

\begin{exercise}
    \(\star\) Un móvil parte de la posición \qty{6}{\m} y se desplaza con MRU en sentido positivo a \qty{12}{\m\per\s}. Escribe su
    ecuación de movimiento y dibuja su gráfica \(x\)-\(t\).
    \begin{solu}
        \(\textcolor{red}{x = 6 + 12t}\) (SI): una recta de pendiente 12 que corta al eje vertical en 6.
    \end{solu}
\end{exercise}

\begin{exercise}
    \(\star\star\) Dos amigos viven en dos pueblos separados \qty{18}{\km} y salen a la vez, uno hacia el otro, por una carretera
    recta. Se encuentran a \qty{7.5}{\km} del primer pueblo al cabo de \qty{25}{\minute}. ¿Qué velocidad lleva cada uno?
    \begin{solu}
        \(v_1 = \dfrac{\num{7.5}}{25/60} = \textcolor{red}{\qty{18}{\km\per\hour}}\) (\qty{5}{\m\per\s}); \(v_2 = \dfrac{\num{10.5}}{25/60} =
        \textcolor{red}{\qty{25.2}{\km\per\hour}}\) (\qty{7}{\m\per\s}).
    \end{solu}
\end{exercise}

\begin{exercise}
    \(\star\star\) Un tren con MRU tarda \qty{15}{\s} en pasar por delante de Luis y \qty{40}{\s} desde que su cabeza entra en un
    túnel de \qty{1000}{\m} hasta que su cola sale de él. ¿Qué longitud tiene el tren? ¿A qué velocidad va?
    \begin{solu}
        \textcolor{blue}{\(L = 15\,v\) y \(L + 1000 = 40\,v\).} \(\textcolor{red}{v = \qty{40}{\m\per\s}}\) (\qty{144}{\km\per\hour});
        \(\textcolor{red}{L = \qty{600}{\m}}\).
    \end{solu}
\end{exercise}

\begin{exercise}
    \(\star\star\) En un videojuego, el coche de Carla pasa por un cruce a \qty{20}{\m\per\s} y, \qty{8}{\s} después, pasa por el
    mismo cruce otro coche a \qty{30}{\m\per\s}, en el mismo sentido. Los dos llevan MRU. a) Representa sus posiciones frente al
    tiempo. b) ¿Cuándo y a qué distancia del cruce la adelanta? c) ¿Qué distancia los separa \qty{10}{\s} después del adelantamiento?
    \begin{solu}
        \textcolor{blue}{\(x_C = 20t\); \(x_2 = 30(t - 8)\).} b) \(\textcolor{red}{t = \qty{24}{\s}}\) desde que Carla pasó el cruce, a
        \(\textcolor{red}{\qty{480}{\m}}\). c) \(10 \cdot (30 - 20) = \textcolor{red}{\qty{100}{\m}}\).
    \end{solu}
\end{exercise}

\subsection{Movimiento rectilíneo uniformemente acelerado}

\begin{exercise}
    \(\star\) Un vehículo arranca con una aceleración de \qty{1.2}{\m\per\s\squared}. Calcula su velocidad \qty{8}{\s} después y la
    distancia que ha recorrido.
    \begin{solu}
        \(v = \textcolor{red}{\qty{9.6}{\m\per\s}}\); \(\Delta x = \num{0.6} \cdot 64 = \textcolor{red}{\qty{38.4}{\m}}\).
    \end{solu}
\end{exercise}

\begin{exercise}
    \(\star\) Un coche que circula a \qty{54}{\km\per\hour} acelera a razón de \qty{2}{\m\per\s\squared} durante \qty{10}{\s}.
    ¿Qué velocidad alcanza y qué espacio recorre?
    \begin{solu}
        \(v = 15 + 20 = \textcolor{red}{\qty{35}{\m\per\s}}\); \(\Delta x = 150 + 100 = \textcolor{red}{\qty{250}{\m}}\).
    \end{solu}
\end{exercise}

\begin{exercise}
    \(\star\) Un coche pasa de 0 a \qty{100}{\km\per\hour} en \qty{8}{\s}. Calcula su aceleración media y la distancia que recorre.
    \begin{solu}
        \(a = \dfrac{\num{27.8}}{8} = \textcolor{red}{\qty{3.47}{\m\per\s\squared}}\); \(\Delta x = \dfrac{\num{27.8}}{2} \cdot 8 =
        \textcolor{red}{\qty{111}{\m}}\).
    \end{solu}
\end{exercise}

\begin{exercise}
    \(\star\star\) Un patinete eléctrico circula durante \qty{6}{\minute} a \qty{18}{\km\per\hour} por un carril bici recto y
    después acelera a \qty{0.5}{\m\per\s\squared} durante \qty{10}{\s}. Halla: a) la distancia total recorrida; b) la rapidez media
    en todo el recorrido.
    \begin{solu}
        a) \(1800 + (50 + 25) = \textcolor{red}{\qty{1875}{\m}}\). b) \(\dfrac{1875}{370} = \textcolor{red}{\qty{5.07}{\m\per\s}}\).
    \end{solu}
\end{exercise}

\begin{exercise}
    \(\star\star\) Un coche circula a \qty{108}{\km\per\hour} y frena con una aceleración de \qty{-6}{\m\per\s\squared}. ¿Qué
    distancia necesita para detenerse? ¿Cuánto tarda?
    \begin{solu}
        \(\Delta x = \dfrac{30^2}{2 \cdot 6} = \textcolor{red}{\qty{75}{\m}}\); \(t = \dfrac{30}{6} = \textcolor{red}{\qty{5}{\s}}\).
    \end{solu}
\end{exercise}

\begin{exercise}
    \(\star\) Una conductora que va a \qty{25}{\m\per\s} quiere detenerse en \qty{4}{\s} y tiene \qty{45}{\m} hasta el coche de
    delante. ¿Le basta esa distancia?
    \begin{solu}
        \(\Delta x = \dfrac{25 + 0}{2} \cdot 4 = \qty{50}{\m}\): \textcolor{red}{no le basta}; necesita al menos \qty{50}{\m}.
    \end{solu}
\end{exercise}

\begin{exercise}
    \(\star\star\) Un motorista que viaja con velocidad constante frena durante \qty{8}{\s} y recorre \qty{120}{\m} hasta pararse.
    Determina: a) su velocidad inicial; b) la aceleración de frenado; c) la aceleración que necesitaría para pararse en la mitad de
    distancia.
    \begin{solu}
        a) \(v_0 = \dfrac{2 \cdot 120}{8} = \textcolor{red}{\qty{30}{\m\per\s}}\); b) \(\textcolor{red}{\qty{-3.75}{\m\per\s\squared}}\);
        c) \(a = -\dfrac{900}{2 \cdot 60} = \textcolor{red}{\qty{-7.5}{\m\per\s\squared}}\).
    \end{solu}
\end{exercise}

\begin{exercise}
    \(\star\star\) Javier y Elena montan en bicicleta en la misma dirección y sentido con MRUA. Al empezar, Javier va a
    \qty{6}{\m\per\s} con una aceleración de \qty{1}{\m\per\s\squared}, y Elena, a \qty{3}{\m\per\s} con \qty{2}{\m\per\s\squared}.
    a) ¿A qué distancia del inicio está cada uno a los \qty{4}{\s}? b) ¿Qué distancia los separa a los \qty{10}{\s}? c) ¿Cuándo y
    dónde adelanta Elena a Javier?
    \begin{solu}
        a) \(\textcolor{red}{x_J = \qty{32}{\m}}\); \(\textcolor{red}{x_E = \qty{28}{\m}}\). b) \(130 - 110 = \textcolor{red}{\qty{20}{\m}}\)
        (va delante Elena). c) \(6t + \num{0.5}\,t^2 = 3t + t^2\): \(\textcolor{red}{t = \qty{6}{\s}}\), a \(\textcolor{red}{\qty{54}{\m}}\).
    \end{solu}
\end{exercise}

\begin{exercise}
    \(\star\star\star\) Un meteorito cae hacia la Tierra a \qty{12}{\km\per\s} constantes. Cuando está a \qty{90}{\km} de altura,
    alguien lo detecta y, \qty{4}{\s} después, despega del suelo, en vertical y desde el reposo, un dron interceptor con una
    aceleración de \qty{40}{\m\per\s\squared}. a) ¿Cuánto tarda en interceptarlo? b) ¿A qué altura? c) ¿Qué aceleración necesitaría
    para interceptarlo a \qty{300}{\m} de altura?
    \begin{solu}
        \textcolor{blue}{\(y_M = 90000 - 12000t\); \(y_D = 20(t - 4)^2\); \(t^2 + 592t - 4484 = 0\).} a) \(\textcolor{red}{t = \qty{7.48}{\s}}\)
        (\qty{3.48}{\s} después de despegar). b) \(\textcolor{red}{y = \qty{242}{\m}}\). c) El meteorito llega a \qty{300}{\m} en
        \(t = \qty{7.475}{\s}\): \(a = \dfrac{2 \cdot 300}{\num{3.475}^2} = \textcolor{red}{\qty{49.7}{\m\per\s\squared}}\).
    \end{solu}
\end{exercise}

\subsection{Caída libre y lanzamiento vertical}

\begin{exercise}
    \(\star\) ¿Cuánto tarda en llegar al suelo una maceta que cae desde \qty{20}{\m} de altura? ¿Con qué velocidad llega? Desprecia el
    rozamiento con el aire.
    \begin{solu}
        \(t = \sqrt{20/\num{4.9}} = \textcolor{red}{\qty{2.02}{\s}}\); \(v = \textcolor{red}{\qty{19.8}{\m\per\s}}\) hacia abajo.
    \end{solu}
\end{exercise}

\begin{exercise}
    \(\star\) Se lanza un cuerpo verticalmente hacia arriba a \qty{72}{\km\per\hour}. a) ¿Qué altura alcanza? b) ¿Cuánto tarda en
    volver al punto de lanzamiento?
    \begin{solu}
        a) \(\dfrac{20^2}{\num{19.6}} = \textcolor{red}{\qty{20.4}{\m}}\); b) \(2 \cdot \dfrac{20}{\num{9.8}} = \textcolor{red}{\qty{4.08}{\s}}\).
    \end{solu}
\end{exercise}

\begin{exercise}
    \(\star\) Una saltadora se deja caer desde una plataforma de \qty{10}{\m} a la piscina. ¿Con qué velocidad llega al agua? ¿Cuánto
    tiempo está en el aire?
    \begin{solu}
        \(v = \sqrt{2 \cdot \num{9.8} \cdot 10} = \textcolor{red}{\qty{14}{\m\per\s}}\); \(t = \textcolor{red}{\qty{1.43}{\s}}\).
    \end{solu}
\end{exercise}

\begin{exercise}
    \(\star\star\) La aceleración de la gravedad en la superficie de cierto planeta es la tercera parte de la terrestre. Si un robot
    se deja caer desde \qty{150}{\m}, ¿con qué velocidad llega a la superficie? ¿Cuánto tarda?
    \begin{solu}
        \(g' = \qty{3.27}{\m\per\s\squared}\); \(v = \sqrt{2 \cdot \num{3.27} \cdot 150} = \textcolor{red}{\qty{31.3}{\m\per\s}}\);
        \(t = \sqrt{300/\num{3.27}} = \textcolor{red}{\qty{9.58}{\s}}\).
    \end{solu}
\end{exercise}

\begin{exercise}
    \(\star\star\) Durante el asedio a un castillo medieval, los defensores lanzan piedras hacia abajo a \qty{10}{\m\per\s} desde
    \qty{30}{\m} de altura. Calcula: a) el tiempo que tarda una piedra en llegar al suelo; b) su velocidad después de recorrer
    \qty{15}{\m}; c) el tiempo que tarda en recorrer el último metro y la velocidad que gana en él.
    \begin{solu}
        a) \(\num{4.9}\,t^2 + 10t - 30 = 0\): \(\textcolor{red}{t = \qty{1.66}{\s}}\). b) \(v = -\sqrt{100 + 2 \cdot \num{9.8} \cdot 15} =
        \textcolor{red}{\qty{-19.8}{\m\per\s}}\). c) \textcolor{blue}{Llega a \qty{1}{\m} del suelo en \(t = \qty{1.618}{\s}\).}
        \(\textcolor{red}{\Delta t = \qty{0.038}{\s}}\); \(\Delta v = \num{9.8} \cdot \num{0.038} = \textcolor{red}{\qty{0.38}{\m\per\s}}\).
    \end{solu}
\end{exercise}

\begin{exercise}
    \(\star\star\) Un tornillo se desprende de un helicóptero que desciende a \qty{8}{\m\per\s} constantes. a) Calcula la velocidad y
    el desplazamiento del tornillo \qty{4}{\s} después. b) Repite el cálculo si el helicóptero asciende a la misma velocidad.
    \begin{solu}
        a) \(v = -8 - \num{39.2} = \textcolor{red}{\qty{-47.2}{\m\per\s}}\); \(\Delta y = -32 - \num{78.4} = \textcolor{red}{\qty{-110.4}{\m}}\).
        b) \(v = 8 - \num{39.2} = \textcolor{red}{\qty{-31.2}{\m\per\s}}\); \(\Delta y = 32 - \num{78.4} = \textcolor{red}{\qty{-46.4}{\m}}\).
    \end{solu}
\end{exercise}

\begin{exercise}
    \(\star\star\) Desde el balcón más alto de la Giralda de Sevilla, a unos \qty{97}{\m} del suelo, se lanza verticalmente hacia
    arriba una rosa a \qty{10}{\m\per\s}. Halla: a) el tiempo hasta la altura máxima y su valor; b) el tiempo que está la rosa en el
    aire y la velocidad con que llega al suelo; c) dibuja las gráficas \(y\)-\(t\), \(v\)-\(t\) y \(a\)-\(t\).
    \begin{solu}
        a) \(\textcolor{red}{t = \qty{1.02}{\s}}\); \(y_{\text{máx}} = 97 + \dfrac{100}{\num{19.6}} = \textcolor{red}{\qty{102.1}{\m}}\).
        b) \(\num{4.9}\,t^2 - 10t - 97 = 0\): \(\textcolor{red}{t = \qty{5.59}{\s}}\); \(v = 10 - \num{9.8} \cdot \num{5.59} =
        \textcolor{red}{\qty{-44.7}{\m\per\s}}\). c) \textcolor{red}{Una parábola cóncava hacia abajo, una recta decreciente y una
        recta horizontal en \(-g\)}.
    \end{solu}
\end{exercise}

\begin{exercise}
    \(\star\star\) Se lanzan verticalmente hacia arriba, desde el mismo punto y con \qty{0.5}{\s} de diferencia, dos pelotas: la
    primera a \qty{5}{\m\per\s} y la segunda a \qty{7}{\m\per\s}. Determina: a) cuándo se cruzan; b) la altura del cruce; c) la
    velocidad de cada pelota en ese instante.
    \begin{solu}
        \textcolor{blue}{\(5t - \num{4.9}\,t^2 = 7(t - \num{0.5}) - \num{4.9}(t - \num{0.5})^2\).} a) \(\textcolor{red}{t = \qty{0.685}{\s}}\)
        tras lanzar la primera. b) \(\textcolor{red}{y = \qty{1.13}{\m}}\). c) \(\textcolor{red}{v_1 = \qty{-1.71}{\m\per\s}}\) (baja) y
        \(\textcolor{red}{v_2 = \qty{5.19}{\m\per\s}}\) (sube).
    \end{solu}
\end{exercise}

\begin{exercise}
    \(\star\star\) Desde \qty{150}{\m} de altura se lanza una moneda hacia arriba a \qty{15}{\m\per\s}. \qty{2}{\s} después se deja
    caer otra desde el mismo punto. a) ¿Cuánto tardan en encontrarse? b) ¿Qué velocidad tiene cada una en el cruce? c) ¿A qué altura
    se cruzan?
    \begin{solu}
        \textcolor{blue}{\(150 + 15t - \num{4.9}\,t^2 = 150 - \num{4.9}(t - 2)^2 \Rightarrow \num{4.6}\,t = \num{19.6}\).}
        a) \(\textcolor{red}{t = \qty{4.26}{\s}}\) tras lanzar la primera. b) \(\textcolor{red}{v_1 = \qty{-26.8}{\m\per\s}}\);
        \(\textcolor{red}{v_2 = \qty{-22.2}{\m\per\s}}\). c) \(\textcolor{red}{y = \qty{125}{\m}}\).
    \end{solu}
\end{exercise}

\subsection{Composición de movimientos}

\begin{exercise}
    \(\star\) Un barco cruza un río de \qty{90}{\m} de ancho a \qty{6}{\m\per\s}, perpendicularmente a una corriente de
    \qty{2.5}{\m\per\s}. ¿Cuánto tarda y a qué distancia río abajo llega a la otra orilla?
    \begin{solu}
        \(t = \dfrac{90}{6} = \textcolor{red}{\qty{15}{\s}}\); \(x = \num{2.5} \cdot 15 = \textcolor{red}{\qty{37.5}{\m}}\).
    \end{solu}
\end{exercise}

\begin{exercise}
    \(\star\star\) Pedro quiere cruzar un río de \qty{2}{\km} de ancho en una barca que avanza a \qty{30}{\km\per\hour} respecto al
    agua, perpendicularmente a la orilla. La corriente es de \qty{6}{\km\per\hour}. Calcula: a) la velocidad de la barca respecto a la
    orilla; b) el tiempo que tarda en cruzar; c) la distancia que recorre río abajo.
    \begin{solu}
        a) \(\sqrt{900 + 36} = \textcolor{red}{\qty{30.6}{\km\per\hour}}\); b) \(\dfrac{2}{30}~\unit{\hour} = \textcolor{red}{\qty{4}{\minute}}\);
        c) \(6 \cdot \dfrac{2}{30} = \textcolor{red}{\qty{0.4}{\km}}\).
    \end{solu}
\end{exercise}

\begin{exercise}
    \(\star\star\) Una persona rema en una canoa a \qty{2.5}{\m\per\s} respecto al agua, formando un ángulo de \(60^\circ\) con la
    dirección de la corriente (río abajo), para cruzar un río de \qty{50}{\m} de ancho cuya corriente es de \qty{1.5}{\m\per\s}.
    a) Escribe el vector velocidad de la canoa respecto a la orilla y su vector posición en función del tiempo. b) ¿Cuánto tarda en
    cruzar y qué distancia recorre río abajo?
    \begin{solu}
        a) \(\vec v = (\num{2.5}\cos 60^\circ + \num{1.5})\,\ivec + \num{2.5}\,\sen 60^\circ\,\jvec =
        \textcolor{red}{(\num{2.75}\,\ivec + \num{2.17}\,\jvec)~\unit{\m\per\s}}\); \(\textcolor{red}{\vec r = \num{2.75}\,t\,\ivec +
        \num{2.17}\,t\,\jvec}\). b) \(t = \dfrac{50}{\num{2.165}} = \textcolor{red}{\qty{23.1}{\s}}\); \(x = \textcolor{red}{\qty{63.5}{\m}}\).
    \end{solu}
\end{exercise}

\begin{exercise}
    \(\star\) En un aeropuerto, una cinta transportadora de \qty{50}{\m} avanza a \qty{0.8}{\m\per\s}. Calcula cuánto tarda en
    recorrerla una persona: a) quieta sobre la cinta; b) andando a \qty{1.2}{\m\per\s} en el sentido de la cinta; c) andando en sentido
    contrario.
    \begin{solu}
        a) \(\dfrac{50}{\num{0.8}} = \textcolor{red}{\qty{62.5}{\s}}\); b) \(\dfrac{50}{2} = \textcolor{red}{\qty{25}{\s}}\);
        c) \(\dfrac{50}{\num{0.4}} = \textcolor{red}{\qty{125}{\s}}\) (avanza al revés que la cinta).
    \end{solu}
\end{exercise}

\subsection{Tiro horizontal}

\begin{exercise}
    \(\star\) Desde una ventana a \qty{8}{\m} de altura sale horizontalmente el agua de una manguera a \qty{10}{\m\per\s}. ¿A qué
    distancia de la pared cae el agua?
    \begin{solu}
        \(t = \sqrt{8/\num{4.9}} = \qty{1.28}{\s}\); \(x = \textcolor{red}{\qty{12.8}{\m}}\).
    \end{solu}
\end{exercise}

\begin{exercise}
    \(\star\) Una pelota sale rodando por el borde de una mesa de \qty{80}{\cm} de altura a \qty{1.5}{\m\per\s}. ¿Cuánto tarda en
    llegar al suelo? ¿Con qué velocidad lo hace? ¿A qué distancia de la mesa cae?
    \begin{solu}
        \(\textcolor{red}{t = \qty{0.404}{\s}}\); \(\textcolor{red}{\vec v = (\num{1.5}\,\ivec - \num{3.96}\,\jvec)~\unit{\m\per\s}}\);
        \(x = \textcolor{red}{\qty{0.61}{\m}}\).
    \end{solu}
\end{exercise}

\begin{exercise}
    \(\star\star\) Javier sube a la azotea de un edificio de \qty{15}{\m} y lanza un balón horizontalmente a \qty{18}{\m\per\s}.
    Despreciando el rozamiento, calcula: a) la distancia al edificio a la que cae; b) el tiempo que tarda; c) su velocidad al llegar
    al suelo.
    \begin{solu}
        a) \(\textcolor{red}{\qty{31.5}{\m}}\); b) \(\textcolor{red}{\qty{1.75}{\s}}\); c) \(\textcolor{red}{\vec v = (18\,\ivec -
        \num{17.1}\,\jvec)~\unit{\m\per\s}}\), de módulo \qty{24.8}{\m\per\s}.
    \end{solu}
\end{exercise}

\begin{exercise}
    \(\star\star\) Desde la cumbre del Mulhacén, Nieves lanza un plato volador horizontalmente a \qty{50}{\m\per\s} hacia un valle
    situado \qty{1000}{\m} más abajo. Suponiendo que nada lo frena, determina: a) cuánto tarda en llegar al valle; b) a qué
    distancia horizontal cae; c) su velocidad y su posición, respecto al punto de lanzamiento, a los \qty{10}{\s}.
    \begin{solu}
        a) \(t = \sqrt{1000/\num{4.9}} = \textcolor{red}{\qty{14.3}{\s}}\); b) \(\textcolor{red}{\qty{714}{\m}}\);
        c) \(\textcolor{red}{\vec v = (50\,\ivec - 98\,\jvec)~\unit{\m\per\s}}\); \(\textcolor{red}{\vec r = (500\,\ivec - 490\,\jvec)~\unit{\m}}\).
    \end{solu}
\end{exercise}

\begin{exercise}
    \(\star\star\) Desde un helicóptero parado a \qty{45}{\m} sobre el mar, un socorrista se lanza horizontalmente a
    \qty{8}{\m\per\s} para rescatar a una persona. Calcula: a) a qué distancia del helicóptero (en horizontal) cae al agua; b) el
    tiempo que está en el aire; c) su velocidad al llegar al agua.
    \begin{solu}
        a) \(\textcolor{red}{\qty{24.2}{\m}}\); b) \(\textcolor{red}{\qty{3.03}{\s}}\); c) \(\textcolor{red}{\vec v = (8\,\ivec -
        \num{29.7}\,\jvec)~\unit{\m\per\s}}\).
    \end{solu}
\end{exercise}

\begin{exercise}
    \(\star\star\) Un avión vuela horizontalmente a \qty{800}{\m} de altura y a \qty{360}{\km\per\hour} y suelta un paquete.
    Determina: a) el tiempo que tarda en llegar al suelo; b) la distancia horizontal que recorre; c) a qué distancia del avión está el
    paquete cuando toca el suelo, si el avión sigue igual.
    \begin{solu}
        a) \(\textcolor{red}{\qty{12.8}{\s}}\); b) \(\textcolor{red}{\qty{1278}{\m}}\); c) \textcolor{red}{A \qty{800}{\m}, justo
        debajo}: los dos avanzan lo mismo en horizontal.
    \end{solu}
\end{exercise}

\subsection{Tiro parabólico}

\begin{exercise}
    \(\star\star\) Jorge salta con una moto de agua que va a \qty{54}{\km\per\hour} por una ola que le hace despegar con un ángulo de
    \(15^\circ\). Calcula: a) la distancia que recorre en el aire; b) el tiempo que está en el aire; c) la posición del punto más alto.
    \begin{solu}
        a) \(\dfrac{225\,\sen 30^\circ}{\num{9.8}} = \textcolor{red}{\qty{11.5}{\m}}\); b) \(\textcolor{red}{\qty{0.79}{\s}}\);
        c) \(\textcolor{red}{(\num{5.74}; \num{0.77})~\unit{\m}}\).
    \end{solu}
\end{exercise}

\begin{exercise}
    \(\star\star\) Desde la azotea de un chalet, a \qty{15}{\m} de altura, se lanza un globo de agua a \qty{25}{\m\per\s} con un
    ángulo de \(30^\circ\) sobre la horizontal. Determina: a) a qué distancia de la base de la casa cae; b) el tiempo que está en el
    aire; c) su velocidad al llegar al suelo; d) su vector posición en \(t = \qty{1}{\s}\).
    \begin{solu}
        \textcolor{blue}{\(x = \num{21.65}\,t\); \(y = 15 + \num{12.5}\,t - \num{4.9}\,t^2\).} b) \(\textcolor{red}{t = \qty{3.44}{\s}}\);
        a) \(\textcolor{red}{x = \qty{74.5}{\m}}\); c) \(\textcolor{red}{\vec v = (\num{21.7}\,\ivec - \num{21.2}\,\jvec)~\unit{\m\per\s}}\),
        de módulo \qty{30.3}{\m\per\s}; d) \(\textcolor{red}{\vec r = (\num{21.7}\,\ivec + \num{22.6}\,\jvec)~\unit{\m}}\).
    \end{solu}
\end{exercise}

\begin{exercise}
    \(\star\star\) Una atleta de triple salto se pregunta qué velocidad necesita en un salto con un ángulo de \(30^\circ\) para
    recorrer \qty{14}{\m} (considera el salto como un tiro parabólico de su centro de masas).
    \begin{solu}
        \(14 = \dfrac{v_0^2\,\sen 60^\circ}{\num{9.8}}\): \(\textcolor{red}{v_0 = \qty{12.6}{\m\per\s}}\).
    \end{solu}
\end{exercise}

\begin{exercise}
    \(\star\star\) Un jugador de baloncesto lanza a canasta desde \qty{7}{\m} de distancia horizontal, a \qty{8.8}{\m\per\s} y con un
    ángulo de \(45^\circ\). El balón sale a \qty{2.1}{\m} del suelo y el aro está a \qty{3.05}{\m}. ¿A qué altura pasa el balón por la
    vertical del aro? ¿Encestará?
    \begin{solu}
        \textcolor{blue}{\(v_{0x} = \qty{6.22}{\m\per\s}\); llega a \(x = 7\) en \(t = \qty{1.12}{\s}\).} \(y = \num{2.1} + \num{6.22} \cdot
        \num{1.12} - \num{4.9} \cdot \num{1.12}^2 = \textcolor{red}{\qty{2.90}{\m}}\): \textcolor{red}{se queda corto}, pasa
        \qty{15}{\cm} por debajo del aro.
    \end{solu}
\end{exercise}

\begin{exercise}
    \(\star\star\star\) Un atleta lanza la jabalina desde \qty{2}{\m} de altura a \qty{28}{\m\per\s}. ¿Con qué ángulo debe lanzarla
    para que caiga a \qty{60}{\m}? (Usa \(\dfrac{1}{\cos^2\alpha} = 1 + \tg^2\alpha\).)
    \begin{solu}
        \textcolor{blue}{Ecuación de la trayectoria con \(y = 0\), \(x = 60\): \(0 = 2 + 60\,\tg\alpha - \num{22.5}(1 + \tg^2\alpha)\),
        es decir, \(\num{22.5}\,\tg^2\alpha - 60\,\tg\alpha + \num{20.5} = 0\).} \(\tg\alpha = \num{2.26}\) o \num{0.402}:
        \(\textcolor{red}{\alpha = 66.2^\circ}\) o \(\textcolor{red}{\alpha = 21.9^\circ}\).
    \end{solu}
\end{exercise}

\begin{exercise}
    \(\star\star\star\) María vive en un cuarto piso (cada piso mide \qty{2.8}{\m}) y lanza con un tirachinas una caja pequeña a
    \qty{20}{\m\per\s} con un ángulo de \(45^\circ\) \textbf{por debajo} de la horizontal hacia Pilar, que está en la calle.
    Determina: a) a qué distancia del edificio debe colocarse Pilar; b) el tiempo que la caja está en el aire y su velocidad al
    llegar al suelo.
    \begin{solu}
        \textcolor{blue}{\(y_0 = \qty{11.2}{\m}\); \(v_{0x} = \qty{14.1}{\m\per\s}\); \(v_{0y} = \qty{-14.1}{\m\per\s}\);
        \(\num{4.9}\,t^2 + \num{14.14}\,t - \num{11.2} = 0\).} b) \(\textcolor{red}{t = \qty{0.65}{\s}}\);
        \(\textcolor{red}{\vec v = (\num{14.1}\,\ivec - \num{20.5}\,\jvec)~\unit{\m\per\s}}\). a) \(x = \textcolor{red}{\qty{9.1}{\m}}\).
    \end{solu}
\end{exercise}

\subsection{Movimiento circular uniforme}

\begin{exercise}
    \(\star\) Una partícula describe una circunferencia de \qty{40}{\cm} de radio a \qty{30}{rpm}. Expresa su periodo, su frecuencia,
    su velocidad angular y su velocidad lineal en el SI.
    \begin{solu}
        \(\textcolor{red}{T = \qty{2}{\s}}\); \(\textcolor{red}{f = \qty{0.5}{\hertz}}\); \(\textcolor{red}{\omega = \pi = \qty{3.14}{\radian\per\s}}\);
        \(\textcolor{red}{v = \qty{1.26}{\m\per\s}}\).
    \end{solu}
\end{exercise}

\begin{exercise}
    \(\star\) Una lavadora centrifuga a \qty{1000}{rpm} y su tambor tiene \qty{50}{\cm} de diámetro. Calcula la velocidad angular, la
    velocidad lineal de la ropa pegada al tambor, el periodo y la frecuencia.
    \begin{solu}
        \(\textcolor{red}{\omega = \qty{105}{\radian\per\s}}\); \(\textcolor{red}{v = \qty{26.2}{\m\per\s}}\); \(\textcolor{red}{T = \qty{0.06}{\s}}\);
        \(\textcolor{red}{f = \qty{16.7}{\hertz}}\).
    \end{solu}
\end{exercise}

\begin{exercise}
    \(\star\) El volante de un triciclo, de \qty{40}{\cm} de diámetro, da 6 vueltas por minuto. Calcula su velocidad angular, la
    velocidad lineal de un punto del borde y su aceleración normal.
    \begin{solu}
        \(\textcolor{red}{\omega = \qty{0.628}{\radian\per\s}}\); \(\textcolor{red}{v = \qty{0.126}{\m\per\s}}\);
        \(\textcolor{red}{a_n = \qty{0.079}{\m\per\s\squared}}\).
    \end{solu}
\end{exercise}

\begin{exercise}
    \(\star\star\) Una mosca se posa en el extremo del aspa de un ventilador de \qty{60}{\cm} de diámetro que gira a \qty{90}{rpm}.
    Determina: a) su velocidad angular y lineal; b) su periodo y su frecuencia; c) su aceleración normal; d) el ángulo y el arco que
    recorre en \qty{2}{\s}.
    \begin{solu}
        a) \(\textcolor{red}{\qty{9.42}{\radian\per\s}}\) y \(\textcolor{red}{\qty{2.83}{\m\per\s}}\); b) \(\textcolor{red}{\qty{0.667}{\s}}\) y
        \(\textcolor{red}{\qty{1.5}{\hertz}}\); c) \(\textcolor{red}{\qty{26.6}{\m\per\s\squared}}\); d) \(\textcolor{red}{\qty{18.8}{\radian}}\)
        y \(\textcolor{red}{\qty{5.65}{\m}}\).
    \end{solu}
\end{exercise}

\begin{exercise}
    \(\star\star\) ¿Cuál es la velocidad angular de cada una de las agujas de un reloj analógico?
    \begin{solu}
        Horaria (una vuelta en \qty{12}{\hour}): \(\textcolor{red}{\qty{1.45e-4}{\radian\per\s}}\); minutero (\qty{1}{\hour}):
        \(\textcolor{red}{\qty{1.75e-3}{\radian\per\s}}\); segundero (\qty{1}{\minute}): \(\textcolor{red}{\qty{0.105}{\radian\per\s}}\).
    \end{solu}
\end{exercise}

\begin{exercise}
    \(\star\star\star\) La Estación Espacial Internacional (ISS) da una vuelta a la Tierra cada \qty{92}{\minute} en una órbita
    prácticamente circular a unos \qty{410}{\km} de altura. Si el radio de la Tierra es de \qty{6370}{\km}, calcula: a) su velocidad
    angular, su periodo y su frecuencia; b) su velocidad lineal; c) su aceleración normal.
    \begin{solu}
        \textcolor{blue}{\(R = \qty{6.78e6}{\m}\); \(T = \qty{5520}{\s}\).} a) \(\textcolor{red}{\omega = \qty{1.14e-3}{\radian\per\s}}\);
        \(\textcolor{red}{f = \qty{1.81e-4}{\hertz}}\). b) \(\textcolor{red}{v = \qty{7.72e3}{\m\per\s}}\) (casi \qty{28000}{\km\per\hour}).
        c) \(\textcolor{red}{a_n = \qty{8.78}{\m\per\s\squared}}\): casi la misma \(g\) que en la superficie.
    \end{solu}
\end{exercise}

\subsection{Movimiento circular uniformemente acelerado}

\begin{exercise}
    \(\star\) Una rueda de un tractor alcanza las 60 vueltas por minuto \qty{15}{\s} después de arrancar desde el reposo. Calcula su
    aceleración angular, el ángulo que describe en esos \qty{15}{\s} y las vueltas que da.
    \begin{solu}
        \(\textcolor{red}{\alpha = \qty{0.419}{\radian\per\s\squared}}\); \(\textcolor{red}{\Delta\varphi = \qty{47.1}{\radian}}\);
        \(\textcolor{red}{\num{7.5}~\text{vueltas}}\).
    \end{solu}
\end{exercise}

\begin{exercise}
    \(\star\star\) Al levantar la aguja de un tocadiscos, el disco, de \qty{30}{\cm} de diámetro, que giraba a \qty{45}{rpm}, se
    detiene en \qty{12}{\s}. Determina: a) la aceleración angular y la velocidad lineal de un punto del borde a los \qty{4}{\s};
    b) las vueltas que da hasta pararse.
    \begin{solu}
        a) \(\textcolor{red}{\alpha = \qty{-0.393}{\radian\per\s\squared}}\); \(\omega(4) = \qty{3.14}{\radian\per\s}\) y
        \(\textcolor{red}{v = \qty{0.47}{\m\per\s}}\). b) \(\Delta\varphi = \qty{28.3}{\radian}\): \(\textcolor{red}{\num{4.5}~\text{vueltas}}\).
    \end{solu}
\end{exercise}

\begin{exercise}
    \(\star\star\) El volante de un coche de carreras gira desde el reposo con una aceleración angular de \qty{2}{\radian\per\s\squared}.
    a) ¿Cuánto tarda en girar \(90^\circ\)? b) ¿Qué velocidad angular tiene entonces?
    \begin{solu}
        a) \(\dfrac{\pi}{2} = t^2\): \(\textcolor{red}{t = \qty{1.25}{\s}}\). b) \(\textcolor{red}{\omega = \qty{2.51}{\radian\per\s}}\).
    \end{solu}
\end{exercise}

\begin{exercise}
    \(\star\star\) Una motocicleta da vueltas a un circuito circular de \qty{150}{\m} de diámetro. Parte del reposo con MCUA y,
    \qty{25}{\s} después, circula a \qty{54}{\km\per\hour}; a partir de ahí sigue con MCU. a) Calcula sus aceleraciones tangencial,
    normal, total y angular justo antes de llegar a los \qty{25}{\s}. b) ¿Qué velocidad angular alcanza? ¿Cuál es entonces su
    periodo?
    \begin{solu}
        a) \(\textcolor{red}{a_t = \qty{0.6}{\m\per\s\squared}}\); \(\textcolor{red}{a_n = \dfrac{225}{75} = \qty{3}{\m\per\s\squared}}\);
        \(\textcolor{red}{a = \qty{3.06}{\m\per\s\squared}}\); \(\textcolor{red}{\alpha = \qty{0.008}{\radian\per\s\squared}}\).
        b) \(\textcolor{red}{\omega = \qty{0.2}{\radian\per\s}}\); \(\textcolor{red}{T = \qty{31.4}{\s}}\).
    \end{solu}
\end{exercise}

\begin{exercise}
    \(\star\star\) Un engranaje gira a \qty{20}{\radian\per\s} y frena con una aceleración angular de \qty{-4}{\radian\per\s\squared}.
    Calcula el tiempo que tarda en detenerse y el ángulo que recorre.
    \begin{solu}
        \(\textcolor{red}{t = \qty{5}{\s}}\); \(\Delta\varphi = \dfrac{20}{2} \cdot 5 = \textcolor{red}{\qty{50}{\radian}}\).
    \end{solu}
\end{exercise}

%========================================================================================
\section{Ejercicios de ampliación}\label{sec:ud9-ampliacion}
%========================================================================================

\begin{exercise}
    \(\star\star\) Esta tabla recoge la velocidad de un coche durante una carrera:
    \begin{center}
        \begin{tabular}{@{}lcccccccc@{}}
            \toprule
            \(\boldsymbol{t}\) \textbf{(s)} & 0 & 5 & 10 & 15 & 20 & 25 & 30 & 35 \\
            \midrule
            \(\boldsymbol{v}\) \textbf{(m/s)} & 4 & 4 & 4 & 12 & 20 & 20 & 20 & 8 \\
            \bottomrule
        \end{tabular}
    \end{center}
    a) Representa la gráfica \(v\)-\(t\). b) Identifica el tipo de movimiento de cada tramo y su aceleración. c) Calcula el espacio
    recorrido en los \qty{35}{\s}.
    \begin{solu}
        b) \textcolor{red}{0-10 s: MRU (\(a = 0\)); 10-20 s: MRUA (\(a = \qty{1.6}{\m\per\s\squared}\)); 20-30 s: MRU; 30-35 s: MRUA de
        frenado (\(a = \qty{-2.4}{\m\per\s\squared}\))}. c) Área bajo la gráfica: \(40 + 120 + 200 + 70 = \textcolor{red}{\qty{430}{\m}}\).
    \end{solu}
\end{exercise}

\begin{exercise}
    \(\star\star\) Arantxa entra corriendo en el andén a \qty{4}{\m\per\s} justo cuando su tren arranca con una aceleración de
    \qty{0.5}{\m\per\s\squared}. Entre ella y la puerta del último vagón hay \qty{15}{\m}. ¿Conseguirá subir al tren?
    \begin{solu}
        \textcolor{blue}{\(4t = 15 + \num{0.25}\,t^2 \Rightarrow t^2 - 16t + 60 = 0\).} \textcolor{red}{Sí}: alcanza la puerta a los
        \(\textcolor{red}{\qty{6}{\s}}\), cuando el tren ha avanzado \qty{9}{\m}. (A los 6 s el tren va a \qty{3}{\m\per\s}, más
        despacio que ella.)
    \end{solu}
\end{exercise}

\begin{exercise}
    \(\star\star\star\) La protagonista de una serie huye por un camino que cruza perpendicularmente una vía de tren. Un tren se
    acerca al cruce a \qty{72}{\km\per\hour} y está a \qty{400}{\m} cuando ella, que está parada a \qty{60}{\m} del cruce, empieza a
    correr hacia él con aceleración constante. a) ¿Qué aceleración necesita para llegar al cruce justo cuando llega el tren? b) ¿Con qué
    velocidad llega?
    \begin{solu}
        \textcolor{blue}{El tren llega en \(t = 400/20 = \qty{20}{\s}\).} a) \(60 = \frac{1}{2}\,a \cdot 400\): \(\textcolor{red}{a =
        \qty{0.3}{\m\per\s\squared}}\). b) \(v = \num{0.3} \cdot 20 = \textcolor{red}{\qty{6}{\m\per\s}}\).
    \end{solu}
\end{exercise}

\begin{exercise}
    \(\star\star\) Una maceta cae desde el poyete de una ventana. Manuel ve que pasa por delante de la suya, que está a \qty{8}{\m} del
    suelo, \qty{1.5}{\s} después de caer. ¿A qué altura estaba el poyete?
    \begin{solu}
        \(h = 8 + \num{4.9} \cdot \num{1.5}^2 = \textcolor{red}{\qty{19.0}{\m}}\).
    \end{solu}
\end{exercise}

\begin{exercise}
    \(\star\star\star\) Desde lo alto de un acantilado se deja caer una piedra, que recorre la última cuarta parte de la altura en
    \qty{1}{\s}. ¿Qué altura tiene el acantilado?
    \begin{solu}
        \textcolor{blue}{Si la caída dura \(T\): \(h = \num{4.9}\,T^2\) y \(\frac{3}{4}h = \num{4.9}(T - 1)^2\), así que \(T - 1 =
        \sqrt{3/4}\,T\) y \(T = \qty{7.46}{\s}\).} \(h = \textcolor{red}{\qty{273}{\m}}\).
    \end{solu}
\end{exercise}

\begin{exercise}
    \(\star\star\star\) Un cohete parte del reposo y sube verticalmente con una aceleración de \qty{2}{\m\per\s\squared}. A los
    \qty{15}{\s} se le acaba el combustible y sigue en caída libre. a) ¿Llegará a \qty{300}{\m} de altura? b) ¿Con qué velocidad
    llega al suelo? c) ¿Cuánto tiempo está en el aire? d) ¿Qué aceleración necesitaría para llegar a \qty{300}{\m}?
    \begin{solu}
        \textcolor{blue}{A los 15 s: \(v = \qty{30}{\m\per\s}\), \(y = \qty{225}{\m}\); después sube \(30^2/\num{19.6} = \qty{45.9}{\m}\) más.}
        a) \textcolor{red}{No: llega a \qty{271}{\m}}. b) \(\sqrt{2 \cdot \num{9.8} \cdot 271} = \textcolor{red}{\qty{72.9}{\m\per\s}}\).
        c) \(15 + \num{3.06} + \num{7.44} = \textcolor{red}{\qty{25.5}{\s}}\). d) \(\num{112.5}\,a + \num{11.48}\,a^2 = 300\):
        \(\textcolor{red}{a = \qty{2.18}{\m\per\s\squared}}\).
    \end{solu}
\end{exercise}

\begin{exercise}
    \(\star\star\) Carmen viaja en un tren que va a \qty{15}{\m\per\s} y lanza su móvil verticalmente hacia arriba a
    \qty{4}{\m\per\s}. a) Para Carmen, ¿cuánto tarda el móvil en volver a sus manos y con qué velocidad lo hace? b) Repite los
    cálculos para una persona que lo observa desde el andén.
    \begin{solu}
        a) \(\textcolor{red}{t = \qty{0.82}{\s}}\); \(\textcolor{red}{\vec v = -4\,\jvec~\unit{\m\per\s}}\). b) \textcolor{red}{El mismo
        tiempo}, pero el móvil describe una parábola: vuelve con \(\textcolor{red}{\vec v = (15\,\ivec - 4\,\jvec)~\unit{\m\per\s}}\) tras
        avanzar \qty{12.2}{\m} con el tren.
    \end{solu}
\end{exercise}

\begin{exercise}
    \(\star\star\star\) Un avión vuela hacia el oeste a \qty{600}{\km\per\hour} respecto al aire, y el viento sopla de norte a sur a
    \qty{100}{\km\per\hour}. Calcula: a) la velocidad del avión respecto al suelo (vector y módulo); b) el ángulo que forma con la
    dirección oeste.
    \begin{solu}
        a) Con \(\ivec\) hacia el este y \(\jvec\) hacia el norte: \(\textcolor{red}{\vec v = (-600\,\ivec - 100\,\jvec)~\unit{\km\per\hour}}\),
        de módulo \(\textcolor{red}{\qty{608}{\km\per\hour}}\). b) \(\textcolor{red}{\num{9.5}^\circ}\) hacia el sur.
    \end{solu}
\end{exercise}

\begin{exercise}
    \(\star\star\) Patricia lanza desde su ventana, a \qty{4}{\m} del suelo, el agua de un cubo horizontalmente a \qty{6}{\m\per\s}.
    a) ¿A qué distancia de la pared cae el agua? b) ¿A qué velocidad habría que lanzarla para que cayera a \qty{2}{\m}?
    \begin{solu}
        a) \(t = \qty{0.904}{\s}\); \(\textcolor{red}{x = \qty{5.4}{\m}}\). b) \(v = \dfrac{2}{\num{0.904}} = \textcolor{red}{\qty{2.2}{\m\per\s}}\).
    \end{solu}
\end{exercise}

\begin{exercise}
    \(\star\star\star\) En un entrenamiento, Carlos lanza un balón desde el suelo a \qty{25}{\m\per\s} con un ángulo de \(45^\circ\).
    En cada bote, el balón sale de nuevo a \(45^\circ\), pero solo alcanza la mitad de la altura máxima anterior. a) ¿Cuál es la
    altura máxima antes y después del primer bote? b) ¿A qué distancia del lanzamiento bota la primera y la segunda vez?
    \begin{solu}
        a) \(\textcolor{red}{\qty{15.9}{\m}}\) y \(\textcolor{red}{\qty{8.0}{\m}}\). b) \textcolor{blue}{La altura máxima es
        proporcional a \(v_0^2\) y el alcance también: si la altura se reduce a la mitad, el alcance también.} Primer bote a
        \(\textcolor{red}{\qty{63.8}{\m}}\); segundo a \(63{,}8 + 31{,}9 = \textcolor{red}{\qty{95.7}{\m}}\).
    \end{solu}
\end{exercise}

\begin{exercise}
    \(\star\star\star\) Desde lo alto de una torre de altura \(h\) se lanzan dos pelotas con la misma rapidez \(v_0\): una hacia
    arriba con un ángulo \(\alpha\) sobre la horizontal y otra hacia abajo con el mismo ángulo. a) Demuestra que las dos llegan al
    suelo con la misma rapidez. b) Calcula esa rapidez en función de \(v_0\), \(g\) y \(h\).
    \begin{solu}
        \textcolor{blue}{Las dos tienen la misma \(v_x = v_0\cos\alpha\), y \(v_y^2 = v_{0y}^2 + 2gh\) con \(v_{0y}^2 = v_0^2\,\sen^2\alpha\)
        en los dos casos (el signo de \(v_{0y}\) desaparece al elevarlo al cuadrado).} a) \textcolor{red}{Llegan con la misma
        rapidez}, aunque en instantes distintos. b) \(\textcolor{red}{v = \sqrt{v_0^2 + 2gh}}\).
    \end{solu}
\end{exercise}

\begin{exercise}
    \(\star\star\star\) Una manzana cuelga a \qty{5}{\m} de altura de un árbol que está a \qty{18}{\m} de Patricia, que dispara una
    flecha desde \qty{1.5}{\m} de altura a \qty{22}{\m\per\s} y con un ángulo de \(12^\circ\). En el instante del disparo, la manzana
    se suelta. a) ¿Alcanza la flecha a la manzana? b) Demuestra que si Patricia apuntara directamente a la manzana
    (\(\alpha = 11{,}0^\circ\)), le daría, cualquiera que fuese la velocidad de la flecha (si llega al árbol antes de que la manzana
    toque el suelo).
    \begin{solu}
        a) \textcolor{blue}{Llega al árbol en \(t = \dfrac{18}{22\cos 12^\circ} = \qty{0.836}{\s}\); \(y_F = \qty{1.90}{\m}\) e
        \(y_M = 5 - \num{4.9} \cdot \num{0.836}^2 = \qty{1.57}{\m}\).} \textcolor{red}{No}: pasa \qty{33}{\cm} por encima.
        b) \textcolor{red}{Las dos caen lo mismo, \(\frac{1}{2}gt^2\), respecto a la recta de puntería}: sin gravedad, la flecha iría en
        línea recta hasta la manzana quieta; con gravedad, las dos «bajan» lo mismo en el mismo tiempo.
    \end{solu}
\end{exercise}

\begin{exercise}
    \(\star\star\) Chloe va en un caballo de un tiovivo, a \qty{2.5}{\m} del eje, con una aceleración normal de \qty{0.5}{\m\per\s\squared}.
    Se le cae un zapato desde \qty{1.2}{\m} de altura. Determina: a) la velocidad lineal, el periodo y la frecuencia del tiovivo;
    b) a qué distancia horizontal del punto en que se soltó cae el zapato.
    \begin{solu}
        a) \(v = \sqrt{a_n R} = \textcolor{red}{\qty{1.12}{\m\per\s}}\); \(\textcolor{red}{T = \qty{14.0}{\s}}\);
        \(\textcolor{red}{f = \qty{0.071}{\hertz}}\). b) El zapato sale con la velocidad tangente: tiro horizontal durante
        \(\sqrt{\num{1.2}/\num{4.9}} = \qty{0.49}{\s}\): \(x = \textcolor{red}{\qty{0.55}{\m}}\).
    \end{solu}
\end{exercise}

%========================================================================================
\section{Olimpiadas de Física}\label{sec:ud9-olimpiadas}
%========================================================================================

\begin{exercise}
    \(\star\star\star\) \textit{(Adaptado de la fase local de Valencia, 2020.)} Un jugador de rugby lanza el balón desde \qty{2}{\m}
    de altura a \qty{12}{\m\per\s} y con un ángulo de \(30^\circ\) sobre la horizontal. Considera el balón como una partícula y
    desprecia el rozamiento. a) ¿A qué distancia llega el balón si el lanzador está quieto? b) Ahora el lanzador corre hacia atrás a
    \qty{0.4}{\m\per\s} mientras lanza (en sentido contrario al del lanzamiento). Otro jugador, en reposo junto a él, empieza a correr
    en ese instante con aceleración constante para coger el balón a \qty{1.5}{\m} de altura. ¿Qué aceleración necesita? ¿Es posible
    en una situación real?
    \begin{solu}
        a) \textcolor{blue}{\(\num{4.9}\,t^2 - 6t - 2 = 0\): \(t = \qty{1.50}{\s}\).} \(x = \num{10.39} \cdot \num{1.50} =
        \textcolor{red}{\qty{15.6}{\m}}\). b) \textcolor{blue}{Respecto al suelo, \(v_{0x} = \num{10.39} - \num{0.4} = \qty{9.99}{\m\per\s}\);
        llega a \qty{1.5}{\m} de altura en \(t = \qty{1.30}{\s}\), a \(x = \qty{13.0}{\m}\).} \(a = \dfrac{2x}{t^2} =
        \textcolor{red}{\qty{15.3}{\m\per\s\squared}}\): \textcolor{red}{imposible para una persona} (un velocista apenas supera
        \qty{5}{\m\per\s\squared} en la salida).
    \end{solu}
\end{exercise}

\begin{exercise}
    \(\star\star\star\) \textit{(Adaptado de la fase local de Navarra, 2021.)} Un dron de mensajería vuela horizontalmente a
    \qty{4}{\m\per\s} y a \qty{40}{\m} de altura hacia la azotea de un edificio de \qty{30}{\m}, donde debe soltar un paquete que caerá
    libremente. a) ¿A qué distancia horizontal de la entrega debe soltarlo? b) Una ráfaga de viento da al paquete una aceleración
    adicional \(\vec a = (-2\,\ivec + \jvec)~\unit{\m\per\s\squared}\) (el dron vuela en el sentido de \(\ivec\)). ¿A qué distancia
    debe soltarlo ahora?
    \begin{solu}
        a) \textcolor{blue}{Cae \qty{10}{\m}: \(t = \sqrt{10/\num{4.9}} = \qty{1.43}{\s}\).} \(x = 4 \cdot \num{1.43} = \textcolor{red}{\qty{5.7}{\m}}\)
        antes del punto de entrega. b) \textcolor{blue}{\(a_y = -\num{9.8} + 1 = \qty{-8.8}{\m\per\s\squared}\): \(t = \sqrt{20/\num{8.8}} =
        \qty{1.51}{\s}\); \(x = 4t - t^2\).} \(x = \textcolor{red}{\qty{3.8}{\m}}\) antes.
    \end{solu}
\end{exercise}

%========================================================================================
\section{Autoevaluación}\label{sec:ud9-autoevaluacion}
%========================================================================================

\begin{exercise}
    Dos móviles están separados \qty{60}{\m} y empiezan a moverse a la vez, uno hacia el otro. El de la izquierda va a
    \qty{5}{\m\per\s} constantes y el de la derecha parte a \qty{1}{\m\per\s} con una aceleración de \qty{2}{\m\per\s\squared}.
    ¿Dónde y cuándo se encuentran?
    \begin{solu}
        \textcolor{blue}{\(5t + t + t^2 = 60\).} \(\textcolor{red}{t = \qty{5.31}{\s}}\), a \(\textcolor{red}{\qty{26.5}{\m}}\) de la posición
        inicial del primero.
    \end{solu}
\end{exercise}

\begin{exercise}
    Desde \qty{30}{\m} de altura se lanza hacia arriba un objeto a \qty{4}{\m\per\s} y, en el mismo instante, desde el suelo se lanza
    otro hacia arriba a \qty{30}{\m\per\s}. Calcula: a) cuándo se encuentran; b) a qué altura; c) si suben o bajan en ese momento.
    \begin{solu}
        a) \(30 + 4t = 30t\): \(\textcolor{red}{t = \qty{1.15}{\s}}\). b) \(\textcolor{red}{y = \qty{28.1}{\m}}\).
        c) \(v_1 = \qty{-7.3}{\m\per\s}\): \textcolor{red}{el primero baja}; \(v_2 = \qty{18.7}{\m\per\s}\): \textcolor{red}{el segundo sube}.
    \end{solu}
\end{exercise}

\begin{exercise}
    Un proyectil se lanza desde el suelo con un ángulo de \(30^\circ\) y una velocidad de \qty{80}{\m\per\s}. Determina: a) el alcance
    y la altura máxima; b) la velocidad del proyectil al llegar al suelo.
    \begin{solu}
        a) \(\textcolor{red}{\qty{566}{\m}}\) y \(\textcolor{red}{\qty{81.6}{\m}}\). b) \textcolor{red}{\qty{80}{\m\per\s}, formando \(30^\circ\)
        por debajo de la horizontal}: \(\vec v = (\num{69.3}\,\ivec - 40\,\jvec)~\unit{\m\per\s}\).
    \end{solu}
\end{exercise}

\begin{exercise}
    Un avión vuela horizontalmente a \qty{1500}{\m} de altura y a \qty{720}{\km\per\hour} y deja caer un objeto. Despreciando el
    rozamiento, calcula: a) dónde toca el suelo; b) la velocidad del objeto a los \qty{2}{\s}.
    \begin{solu}
        a) \(t = \qty{17.5}{\s}\): a \(\textcolor{red}{\qty{3.50}{\km}}\) del punto en que se soltó (en horizontal).
        b) \(\textcolor{red}{\vec v = (200\,\ivec - \num{19.6}\,\jvec)~\unit{\m\per\s}}\), de módulo \qty{201}{\m\per\s}.
    \end{solu}
\end{exercise}

\begin{exercise}
    Una bicicleta recorre \qty{12}{\km} en \qty{30}{\minute} con MRU. Si el radio de sus ruedas es de \qty{35}{\cm}, calcula: a) el
    número de vueltas que dan las ruedas; b) su velocidad angular y la velocidad lineal de un punto de la cubierta.
    \begin{solu}
        a) \(\dfrac{12000}{2\pi \cdot \num{0.35}} = \textcolor{red}{\num{5457}~\text{vueltas}}\). b) \(v = \textcolor{red}{\qty{6.67}{\m\per\s}}\)
        (respecto al eje de la rueda); \(\omega = \dfrac{\num{6.67}}{\num{0.35}} = \textcolor{red}{\qty{19.0}{\radian\per\s}}\).
    \end{solu}
\end{exercise}
